\chapter{Conclusion and Future Scope}
\label{ch:conclusion}

\section{Key Findings}

This project compared conventional, single-agent and multi-agent signal controllers on real
traffic flows and built \ugmarl{}, a multi-agent controller that estimates its own uncertainty and
lets Max-Pressure make the decision whenever that uncertainty is high.
\begin{itemize}
\item \textbf{H1, performance: supported.} On the test flow \ugmarl{} averages
\res{tt}{ugmarl}{clean}\,s of travel time, clearly below Max-Pressure
(\res{tt}{max_pressure}{clean}\,s) and IDQN (\res{tt}{idqn}{clean}\,s), and
\res{regret}{ugmarl}{clean}\,s behind the single shared DQN.
\item \textbf{H2, safety: supported}, with one of its three parts undecided. At detector
noise 0.2 the gate holds travel time at \res{tt}{ugmarl}{noise0.2}\,s, where the same model
without it reaches \res{tt}{ugmarl-nogate}{noise0.2}\,s, and keeps the maximum wait at
\res{maxw}{ugmarl}{noise0.2} instead of \res{maxw}{ugmarl-nogate}{noise0.2}\,s. On clean data
the gate makes no measurable difference.
\item \textbf{H3, coordination: not supported.} Neighbour attention did not reduce travel
time, and the model without it does as well with 63\,\% fewer parameters.
\item \textbf{Reliability across conditions.} Each baseline has a condition where it does
badly: the learned ones under detector noise, Max-Pressure on clean data and in the surge.
The largest regret of \ugmarl{} over the six conditions is \res{regret}{ugmarl}{worst}\,s,
against \res{regret}{max_pressure}{worst}--\res{regret}{dqn}{worst}\,s for the baselines, and
on an unseen city it performs at the level of DQN.
\item \textbf{What the uncertainty measures.} $\sigma$ rises for inputs that look unlike the
training data. This is why the gate helps under detector noise. It is also why the gate fires
during a demand surge, where it makes the longest wait worse.
\end{itemize}

\section{Limitations}

\textbf{Evidence.} The study rests on five seeds and one test flow per city. The confidence
intervals capture the randomness of training and say nothing about how demand varies from day
to day. Differences of a few seconds, such as the gap to DQN or the threshold sweep, are close
to what the study can resolve. Because the validation flow was used again and again to pick
checkpoints, the architecture and the threshold, the validation numbers are optimistic. The
cause of the gap to DQN has not been isolated.

\textbf{Method.} Three heads on a shared encoder produce the uncertainty estimate. The signal is
coarse: it cannot tell faulty detectors from heavy demand, and it barely correlates with the
learning error. The gated policy comes with no stability guarantee, and the fallback puts no
bound on how long one vehicle waits. Vehicles that never get into the network are missing from
the waiting-time metrics. Only Gaussian noise and randomly lost readings were tested, not
detectors that are stuck, biased or delayed.

\textbf{Scope.} Both cities are regular grids with a single vehicle class and no pedestrians,
buses or incidents, and both are simulated in CityFlow. With a 15\,s decision interval, the
absolute numbers cannot be compared with studies that use 10\,s. Only two rule baselines
were run.

\textbf{Ethics and deployment.} The data comes without a stated licence, so it is used only for
coursework and is not redistributed. The trips are anonymous and carry no identifiers, so no
personal data is processed. Fairness between road users was examined through waiting per
approach. The learned controllers share waiting evenly across the four approaches (22--29\,\%
each), and the longest unbroken queue on one approach lasts 849\,s under \ugmarl{} against
1275\,s under Max-Pressure. Individual vehicles can still wait a very long time when demand is
heavy. The controller has only been run in simulation and has no safety certification, and the
gap between the simulator and real roads is unknown. Putting it on a real road would need fixed
limits on red time, pedestrian phases and supervision by traffic engineers.

\section{Future Research Directions}

The results suggest four next steps. (1) A fallback with a hard cap on the red time of each
movement would bound the worst case of a single vehicle by design. (2) The gate needs a signal
that tells unreliable inputs apart from heavy but correctly measured demand, for instance by
adding a queue-based trigger or randomised prior functions to the ensemble. (3) Filtering the detector
readings before they reach the fallback, since without filtering none of the controllers here
beats Max-Pressure under noise. (4) Tuning the learning settings of \ugmarl{} on their own and
running more seeds, to find out whether the gap to DQN comes from the architecture. Tests on
irregular networks and with realistic detector faults would show how far these findings hold.
